\section{Proposal I: carry the height that the consumer uses}

\subsection{The current interface is stronger than the final consumer needs}

The readable proof defines
\begin{equation}
 \operatorname{HasAbsUpToNormPower}(c,f)
 \iff
 \exists C\in\R,\ e\in\Z,\quad
 |c(\vartheta)|^2=C\,N(Q_\vartheta)^e f(\vartheta)^2
 \quad\text{for every }\vartheta.
\end{equation}
Here \(c\) is an actual complex scalar used in exchange or integrality
identities. The function \(f\) records explicit products of Hecke-value ratios.
The definition permits \(C=0\), including a scalar family that vanishes
identically. See
\src{New/Determinant/Chain/Moduli}{\code{Moduli.lean}}.

The proof propagates this description through products, inverses, turns,
triangular chains, and auxiliary forms. It then extracts a global integer
exponent, proves positivity and trivial-twist normalization of the ratio
products, and only at the end applies a height bound.
The final lower-volume argument consumes neither the exponent nor the
explicit list of ratios. It consumes a bound on a family of positive-real
moduli embedded into \(\C\).

This suggests changing the interface at the point where each local
arithmetic identity is first established. The local identity still proves
the cost of one crossing; the global induction carries only additive costs.
The complex scalar itself remains available in all exact map identities.

\subsection{The prior compact implementation and the additional contribution}

The compact edition already implements the essential bookkeeping idea in
\compact{New/Exchange/CostOfSquares}{\code{CostOfSquares.lean}}:
its \code{CostsSquares n c} bounds the height of
\(\norm{c(\theta_\iota)/c(1)}^{-2}\) by \(2ntx^{999}\).
The compact
\compact{New/Exchange/TriangleNorms}{\code{TriangleNorms.lean}} and
\compact{New/Rows/RowVolume}{\code{RowVolume.lean}} propagate these costs.
The compact triangle already uses an existential endomorphism, a mixed
ambient Gram identity, and its determinant formula, rather than exposing
a named global operator product. That part of the architecture also
receives explicit credit here.
The independently pinned compact revision used here is
\code{72fead102f663d5c7a407964a13dc0cf9d977207}.

The proposed migration adapts both of these compact-source ideas.
The additions and source-specific work in this report are:
\begin{enumerate}
\item an exact, proved reverse-power rule for its actual \code{HeightLe},
so squared estimates recover the current unsquared endpoint with no loss;
\item arbitrary additive real budgets, with the exact double sum over
crossings, so later improvements to individual
local costs are not forced through a single uniform integer square count;
\item a guarded certificate calculus and an explicit migration of both
readable consumers, preserving their nonvanishing, boundary normalization,
and main-term contracts.
\end{enumerate}
These are not among the 18 existing proposal bundles audited in
Section~\ref{sec:exclusion}. This is a nonduplication claim about the specified
corpus, not a claim of priority for height homogeneity in general or
for the compact determinant architecture.

\section{Exact power descent for the actual height}
\label{sec:root}

\subsection{Definition and the missing reverse direction}

The index \(\iota\) consists of a rational prime \(p_\iota\) and a field
isomorphism from \(\C\) to the chosen \(p_\iota\)-adic algebraic closure.
The source's definition
\src{New/Reduction/Lattices/Heights}{\code{HeightLe}} says that
\(\HLe(C,y)\) holds if there are
\[
 H_0\in\C^\times,\qquad \alpha\in\R,\qquad
 h:\N\longrightarrow\R_{\ge0}
\]
such that
\begin{align}
 |\sigma(H_0)|&\le e^\alpha
       &&\text{for every }\sigma\in\operatorname{Aut}(\C), \label{eq:arch}\\
 |\iota(H_0y_\iota)|_{p_\iota}&\le
       p_\iota^{h(p_\iota)}
       &&\text{for every }\iota, \label{eq:local}\\
 \alpha+\sum_{p\in T}h(p)\log p&\le C
       &&\text{for every finite set \(T\) of primes}. \label{eq:budget}
\end{align}
There is no integrality condition on \(H_0\), no fixed field containing it,
and no finite-support condition on \(h\). Both \(\alpha\) and \(C\) may be
negative. These details make the following rule available without any
rounding or extension-degree penalty.

\begin{theorem}[Exact positive-power homogeneity]
\label{thm:root}
For every integer \(n>0\), every \(C\in\R\), and every family
\(y_\iota\in\C\),
\[
 \HLe(C,y^n)\quad\Longleftrightarrow\quad
 \HLe(C/n,y).
\]
Equivalently, \(\HLe(nC,y^n)\Longleftrightarrow\HLe(C,y)\).
No nonvanishing assumption on \(y_\iota\) is necessary.
\end{theorem}

\begin{proof}
The forward-power implication is the existing \code{HeightLe.pow}.
For the other direction, let \(H_0,\alpha,h\) witness \(\HLe(C,y^n)\).
Choose \emph{one} \(G\in\C\) with \(G^n=H_0\); then \(G\ne0\).
Use the witness
\[
 G,\qquad \alpha/n,\qquad p\longmapsto h(p)/n.
\]
For every complex automorphism,
\[
 |\sigma(G)|^n=|\sigma(H_0)|
       \le e^\alpha=(e^{\alpha/n})^n.
\]
For every local isomorphism,
\[
 |\iota(Gy_\iota)|_{p_\iota}^{\,n}
 =|\iota(H_0y_\iota^n)|_{p_\iota}
 \le p_\iota^{h(p_\iota)}
 =\bigl(p_\iota^{h(p_\iota)/n}\bigr)^n.
\]
The map \(u\mapsto u^n\) is increasing on nonnegative real numbers and
injective there, so both desired norm inequalities follow. The exponents
remain nonnegative. Finally,
\[
 \alpha/n+\sum_{p\in T}(h(p)/n)\log p
 =\frac{\alpha+\sum_{p\in T}h(p)\log p}{n}\le C/n.
\]
This uses neither logarithms of \(y_\iota\) nor a different root at each prime.
\end{proof}

In particular,
\begin{equation}
 \boxed{\ \HLe(2B,y^2)\ \Longleftrightarrow\ \HLe(B,y).\ }
 \label{eq:square}
\end{equation}
It is this exact equivalence that makes square-cost propagation compatible
with the readable proof's existing unsquared determinant bound.

\subsection{What has been formalized}

\code{lean/HeightPower.lean} reproduces the common-multiplier definition,
abstracting only the dependent local normed fields and their homomorphisms
from \(\C\). Setting the index to \code{New.PrimePadicIso}, the local field
to its \code{PadicAlgCl}, and the map to \code{iso.toRingHom} gives the
source's definition. The proof constructs \(G\) by complex exponentiation
and proves the algebraic identity \(G^n=H_0\) before applying any
automorphism. It does not assume that arbitrary complex automorphisms
commute with a chosen branch of analytic exponentiation.

The checked file includes \code{pow}, \code{of_pow}, \code{pow_iff},
\code{scaled_pow_iff}, and \code{sq_iff}. The archive also supplies
\code{integration/HeightPowerIntegration.lean}, which uses the original
source names. Its full-project import has not been compiled here;
this distinction is recorded in the file itself.

\section{A scalar-cost calculus with the necessary guards}
\label{sec:cost}

\subsection{Definition}

Let \(\theta:\iota\mapsto\theta_\iota\) be one family satisfying the readable
proof's \code{HasSmallRatios}. For a scalar family \(c(\vartheta)\), define
the nonnegative real modulus family, embedded into \(\C\), by
\[
 r_{c,\theta}(\iota)
   =\bigl(\norm{c(1)/c(\theta_\iota)}:\C\bigr).
\]
Equivalently, its square is
\((\norm{c(\theta_\iota)/c(1)}^2:\C)^{-1}\).
The algebraic equality also respects Lean's totalized division convention.
Its interpretation as a ratio of nonzero volumes will require explicit
nonvanishing later.

\begin{definition}
For \(B\in\R\), define the guarded squared-cost assertion
\begin{equation}
 \SC_B(c):\quad
 c(1)\ne0\ \Longrightarrow\
 \forall\theta,\ 
 \operatorname{HasSmallRatios}(\theta)\Longrightarrow
 \HLe(B,r_{c,\theta}^{\,2}).
\end{equation}
\end{definition}

The guard is useful because local scalar identities may initially permit
the zero family. It does not assert that \(c\) is nonzero. Actual
nonvanishing is retained as a separate fact where a map or determinant
argument needs it. An unguarded variant, as in the compact source, is also
possible, but then its zero cases must be handled explicitly.

\begin{proposition}[Additive propagation]
If \(\SC_B(c)\) and \(\SC_{B'}(c')\), then
\(\SC_{B+B'}(cc')\). A nonzero constant scalar has cost zero.
Finite products therefore have the sum of the individual costs.
\end{proposition}

\begin{proof}
If \(c(1)c'(1)\ne0\), both base values are nonzero. For the same family
\(\theta\),
\[
 r_{cc',\theta}^{\,2}
   =r_{c,\theta}^{\,2}r_{c',\theta}^{\,2}.
\]
Apply the existing product rule for \code{HeightLe}. Constants cancel
from normalized ratios. Induction proves the finite-product statement.
\end{proof}

The definition allows signed real budgets. The local crossing costs below
are positive, but the single-lift height contributions contain a negative
main term. The implementation should not insert unnecessary
nonnegativity assumptions into the abstract product calculus.

\subsection{Local costs and conductor powers}

A norm-power-only identity
\[
 |c(\vartheta)|^2=C\,N(Q_\vartheta)^e
\]
with \(c(1)\ne0\) makes its normalized squared modulus a power of
\(N(Q_\vartheta)\). At the local index \(\iota\), the twist conductor is
prime to \(p_\iota\). Consequently that power is a local unit and costs
zero. The same holds for its reciprocal, proved separately.
The source's \code{HeightLe.rpow_mul} already implements a more general
rational-power invariance.

For a ramified critical pair, the local exchange scalar has squared modulus
equal to a fixed constant times an \(L\)-value product. The existing identity
in \src{New/Determinant/Chain/Moduli}{\code{CriticalPair.lValueProd_mul_inv_eq}}
expresses its normalized square as a conductor power times the square of a
product of critical ratios. Each inverse critical ratio costs at most
\(x_1\), and there are at most \(t\) of them. Thus one crossing contributes
\begin{equation}
 \SC_{\,2tx_1}(c_{\mathrm{crossing}}).
 \label{eq:onecrossing}
\end{equation}
One can prove this directly from
\code{lValueProd_eq_mul_prod_ratio_sq}, avoiding a permanent global
\code{lRatioProd} interface. The underlying \(L\)-value identity and
ramification hypotheses remain necessary.

\subsection{Four rules that must not be silently strengthened}

\paragraph{No generic inverse rule.}
An upper bound for \(\HLe(B,y)\) does not give the same bound for
\(\HLe(B,y^{-1})\). The use of \code{HasAbsUpToNormPower.inv} in the
chain's turn argument concerns a norm-power-only scalar. Replace that
particular use by its independently proved reciprocal zero-cost statement.
Do not add an unrestricted \(\SC_B(c)\Rightarrow\SC_B(c^{-1})\).

\paragraph{Modulus is not the complex scalar.}
Equation \eqref{eq:square} produces a bound for
\((\norm{c(1)/c(\theta_\iota)}:\C)\), not for the complex quotient
\(c(1)/c(\theta_\iota)\).
For example,
\[
 z=\frac{2+i}{2-i}
\]
has complex absolute value one, but nonzero opposite valuations at the
two primes of \(\mathbb Q(i)\) above five. Replacing the actual scalar
by its complex modulus inside an arbitrary \(p\)-adic isomorphism is
invalid. Only the modulus families already used by the determinant and
auxiliary-form height consumers are changed.

\paragraph{Costs do not prove nonvanishing.}
The source proves chain injectivity independently of modulus bookkeeping.
Retain that induction. If the final Gram scalar at the trivial twist is
zero, injectivity forces the target space to be zero. In that branch,
choose scalar \(1\) and use the empty determinant convention. In the
nonzero target branch, the exact Gram identity and injectivity provide
the required nonzero scalar values.
Auxiliary-form scalars already constructed in \(\C^\times\) should retain
that type or the equivalent explicit hypotheses.

\paragraph{One globally selected family.}
All factors must use the same \(\theta\). The eventual regular family
satisfies all the single-lift and small-ratio bounds simultaneously.
Selecting a favorable twist independently for each factor would not
justify a product height bound or the common lattice comparison.
